\section{The rate forced by a packet computation}\label{sec:proof}

\begin{theorem}[Theorem A, quantitative packet transfer]\label{thm:main}
\formal{TapeStorageRates/Main.lean}{TapeStorageRates.rate_witnesses_of_growing_widths,TapeStorageRates.rate_witnesses_of_near_full,TapeStorageRates.rate_witnesses_of_full}
Fix a deterministic machine $M$ with $k$ work tapes and a time coefficient $c>0$.
Suppose that for every $n\ge1$ and $\eta>0$ there are $d,b\ge n$ for which
\[
  G_{d,b}\ne\varnothing,\qquad
  \log_2|G_{d,b}|\ge(1-\eta)d^2b .
\]
There is a constant $K=K(M,c)>0$ such that for every sufficiently large
integer $\Delta$ some admissible binary family at cap $\Delta$ satisfies
\[
  \liminf_{i\to\infty}\frac{\log_2|C_i|}{|V_i|}
     \ \ge\ 1-K f_{4k^2+1}(\Delta).
\]
Consequently $R(\Delta)\ge1-K f_{4k^2+1}(\Delta)$ for all sufficiently
large $\Delta$, and $R(\Delta)\to1$.
\end{theorem}
\begin{outline}
Fix $s$ bits per packet. On a sufficiently large successful instance,
conditioning on the incomplete ends of entries, the packet schedule and
the revealed steps leaves a code with a deficit of $O_{M,c}(1)$ bits
per original packet. The trimming estimate below retains at least half
the packets at the chosen cap. Expanding those packets into bits gives
a rate loss of $O_{M,c}(1/s)$. The required degree cap $\Delta_s$ satisfies
$\log_2\log_2\Delta_s=O_{M,c}(s^p\log_2(s+1))$. Inverting it gives
the claimed rate. The following two subsections establish the finite
estimates and make the order of the choices explicit.
\end{outline}

Both thresholds $d,b\ge n$ apply to the same success set. In particular,
the hypothesis holds if $M$ succeeds on $2^{L-o(L)}$ inputs along a sequence
of canonical dimensions tending to infinity. A machine satisfying (T)
succeeds on all $2^L$ inputs at those dimensions. The family in the
conclusion may depend on $\Delta$. Once the cap is fixed, its length
and its minimum paired distance grow together.

\subsection{From successful inputs to a common packet schedule}

Fix a packet size $s\ge1$. Use the hypothesis of \cref{thm:main} with $\eta=1/(2s)$.
We can choose $d$ as large as needed with $b\ge s$ and a nonempty success
set $G=G_{d,b}$ of deficit $h=L-\log_2|G|\le L/(2s)$. Write $b=as+r$
with $a\ge1$ and $0\le r<s$. Each matrix entry consists of $a$ complete
$s$-bit packets followed by $r$ tail bits. There are $n=d^2a$ complete
packets. The packet at $(i,j,\ell)$ is demanded at $(j,i,\ell)$ for $0\le\ell<a$. Since $L=sn+rd^2\le2sn$, we have $h\le n$.

Partition $G$ by the entire string of $rd^2$ tail bits. Some fixed
assignment $z$ leaves a set $F_z$ of successful packet arrays with
\[
  |F_z|\ge |G|\,2^{-rd^2}=2^{sn-h}.
\]
Thus fixing the tails preserves the bound $h$ on the absolute deficit,
with $sn$ as the new ambient bit count. Throughout this computation the
tail of an entry is the fixed string specified by its address in $z$.

The auxiliary input tape uses one cell per packet. The last cell of each
entry holds its $s$ payload bits followed by its fixed tail, so its length
in the original input is $s+r$. All other in-range cells have length $s$.
Each cell outside the input represents $s$ end markers, and each work cell
holds $s$ ordinary work symbols. Starting from any offset in a cell, $s$
source transitions can visit only that cell and its two neighbours.

One batch evaluates these transitions in eleven packet ticks. The input
and work heads move one cell left and read the left, centre and right
cells in succession. A stationary tick evaluates the batch on these
windows. The heads traverse the windows in reverse to write back the
three work cells, return to the old centres, and move by at most one cell
to the new centres. Other work writes preserve the scanned symbols. The final stationary tick may append one packet.
At the start of a batch the pending output is shorter than its next
cell. The batch adds at most $s$ raw output bits, so it completes at
most one cell. On completion, emit its first $s$ bits, pass over its
tail, and keep the suffix for the next cell. Every successful run emits
exactly the $n$ payload packets in transposed order. The public cell
lengths and tail values belong to this auxiliary computation.
\Cref{sec:simulation} will return to ordinary bit tapes.

Use $R=\lfloor\lfloor c(L+1)\rfloor/s\rfloor+1$ batches on every input,
keeping the source configuration fixed after halt. For
$J=2\lceil c\rceil+1$, the bounds $R\le J(n+1)$ and $T=11R\le11J(n+1)$
are uniform in $s,d,b$. Within each batch the varying part of the
schedule consists of the $k+1$ displacements to the new centres and the
packet-append flag. These give at most $(2\cdot3^{k+1})^R$ schedule
words. Hence for an integer $A_0=A_0(k,c)$ some fixed schedule has a
class $W$ satisfying
\[
  \log_2|W|\ \ge\ sn-h-A_0 n\ \ge\ sn-(A_0+1)n.
\]
The compiler further records the finite control, offsets, windows,
pending output, work writes and any output packet. They fit in
$\beta_s=O_M(s)$ bits per tick. On the fixed schedule the input address,
cell type, tick within the batch, and count of previously appended
packets are public. The record at a tick is determined by its input
packet, the preceding control record, and the records from the last
visits to the scanned work cells. Initial values and exterior input
symbols are fixed. This locality and the output at each designated
append tick are proved by \lean{TapeStorageRates.PaidClasses.compiler}.

\subsection{Extracting and separating a code}

Call $W$ a \emph{frozen packet class} on $n=d^2a$ packets if its nonempty
set of packet arrays has common unit walks for the input head, the $k$
work heads and the output position, and its tick records obey the
locality rule above. It also has a designated tick before the deadline
for each demanded packet. At that tick the output position is the
packet's transposed address and the record determines the packet on
$W$. In the compiled class this is its append tick.

\begin{proposition}[Retention at a chosen degree threshold]\label{prop:packet-retention}
\formal{TapeStorageRates/PacketFamily.lean}{TapeStorageRates.packet_code_at_threshold}
Fix integers $k,A\ge0$, $s,\beta\ge1$ and a time coefficient $c_0\ge1$.
Let $D=\operatorname{genieDepth}(k,\beta,c_0,1)$ be the depth bound
of the bounded-stack walk theorem, and put
$B=\sum_{j=0}^{D}(k+1)^j$. Fix a degree threshold $\theta\ge0$, a
radius $u\ge0$, both integers, and a real $\alpha$ with
$0<\alpha<1-B/(\theta+1)$.
For all sufficiently large $d$, every $a\ge1$ and every frozen class $W$
on $n=d^2a$ packets with $\beta$-bit records, at most $c_0(n+1)$ ticks
and $\log_2|W|\ge sn-An$, there is an integer $m$ with $\alpha n\le m\le n$ and a binary
storage code on $sm$ coordinates.
Its degree is at most $s\max(B,\theta)$, every paired distance exceeds
$u$, and its rate is at least $1-(A+1)/(\alpha s)$.
\end{proposition}

\begin{outline}
Point each tick back to its control and work predecessors. The walk
theorem supplies a set $\Gamma$ of ticks with $\beta|\Gamma|\le n$.
On a dependency path starting at an output tick and avoiding $\Gamma$
after its first vertex, the depth bound is $D$.
Fix a most frequent joint value for the records at $\Gamma$. This
costs at most $n$ bits in log cardinality. For a source packet $q$, let
$N(q)$ contain the input addresses reached by these paths from its
designated output tick. Two arrays in the resulting fibre that agree on
$N(q)$ have the same output record there. Indeed the revealed records
already agree, and the locality rule gives agreement at the remaining
ticks by induction in time. Thus $N(q)$ recovers $q$. This application
uses the whole-step walk cut of \cite[Corollary~O(a)]{PaperFourteen} and
the common-schedule recovery argument of \cite[proof of Corollary~6.4]{PaperEight}.

Each tick has at most $k+1$ predecessors. Consequently $|N(q)|\le B$,
and the sum of all input-port degrees is at most $Bn$. Remove both ports
of each coordinate whose input degree exceeds $\theta$. There are at
most $Bn/(\theta+1)$ such coordinates. In the induced graph the maximum
degree is at most $\max(B,\theta)$. The paired-path packing theorem bounds
the number of remaining matched pairs at distance at most $u$ by
$O_{k,c_0,D,\theta,u}(n/\log_2 d)$, uniformly in $a$. At bounded degree, a positive density of short paths would yield many
vertex-disjoint ones. Following their dependency edges forwards and
backwards through the walks would transport the corresponding input
coordinates to their transposed positions. The row-entropy transport
bound limits the number that can be served in $O(n)$ ticks. This is
where the transposed addresses enter the estimate. The formal packing
bound uses the walks, the cut and the positions of the transposed demands
\cite[Theorem~3.3 and Remark~3.9]{PaperFive}. At block scale one, the Lean estimate
\lean{MergeDepth.skel_paired_packing} underlying
\cite[Proposition~4.11(a)]{PaperEleven} also permits $k=0$. The strict inequality defining $\alpha$
leaves room to remove these short pairs and retain $m\ge\alpha n$.
Every surviving path was already a path before deletion.

Fix the omitted packets to common values when projecting the code through
the two deletions. The total cost is at most $s(n-m)$ bits, one $s$-bit
payment for each omitted coordinate. Its former occurrences in the
recovery rules are now constants \cite[Corollary~3.8]{PaperFive}. The
projected code has log cardinality at least $sm-(A+1)n$.
Replace every surviving packet by its $s$ bits and each incidence by
all $s^2$ incidences between the corresponding copies. The bits of the
helper packets determine each demanded bit. Expansion preserves code
cardinality, multiplies the degree bound by $s$, and sends every paired
path back to a packet walk of the same length when its bit indices are
forgotten \cite[Lemma~8.1]{PaperThirteen}. Dividing the cardinality
bound by $sm$ gives the rate in the statement.
\end{outline}

\begin{outline}[Completion of the proof of Theorem A]
Apply \cref{prop:packet-retention} with $c_0=11J$, $\beta=\beta_s$ and
$A=A_0+1$ from the schedule count. Taking $\theta=2B$ permits
$\alpha=1/2$, since $B/(2B+1)<1/2$. Thus the cap is $\Delta_s=2sB$
and the rate loss is at most $2(A+1)/s$. At cap $10sB$, taking $\theta=10B$ permits
a retained fraction of at least $9/10$.
For each successive radius $u$, take a qualifying $d,b$ large enough
for the proposition and with $d\ge2(u+1)$. Their lengths tend to infinity, and
the $u$th code has every paired distance greater than $u$.

It remains to fit $\Delta_s$ below a prescribed $\Delta$.
Write $x=4\beta_s c_0=O_{M,c}(s)$. In the walk cut the full depth is
$D=K_0^{P_0}\lceil x\rceil$, with its radix $K_0$ and digit count $P_0$
taken at cut fraction $1/x$ \cite[Definition~7.1]{PaperFourteen}.
For fixed $k$ the parameter bounds give
$\log_2 K_0=O_k(x)$ and
$P_0=O_k(x^{4k^2}\log_2(x+1))$.
Since $B\le(k+2)^{D+1}$, these inequalities yield
\[
   \log_2\log_2\Delta_s
      =O_{M,c}\!\left(s^{\,4k^2+1}\log_2(s+1)\right).
\]
With $y=\log_2\log_2\Delta$ and $p=4k^2+1$, choose $s$ as the
integer part of a sufficiently small fixed multiple of
$(y/\log_2 y)^{1/p}$. For all sufficiently large $\Delta$ this gives
$\Delta_s\le\Delta$ and $1/s=O_{M,c}(f_p(\Delta))$.
The family at that $s$ proves the limiting lower rate bound.
The supremum $R(\Delta)$ is at least this family's limiting upper rate,
and every binary code has rate at most one. This proves the convergence
asserted in \cref{thm:main}.
\end{outline}

